\section{Conclusion}

The EDAP model formalizes an endogenous mechanism of civilizational stagnation: intraspecific economic cannibalism dissipates adaptive potential, trapping societies in a feudal equilibrium from which interstellar expansion is impossible. This is the \emph{Great Plateau}---an explanation for the observed absence of galactic colonization. The post-Roman Dark Ages exemplify this mechanism: a documented case of technological regression, institutional collapse, and half-millennium stagnation consistent with the model's feudal trap attractor.

The model further suggests, within its existing mathematical structure, the conditions for a transition from plateau to \emph{filter}. The technofeudal threshold $T_{\text{sing}}$ marks the point beyond which automation may decouple elite welfare from population size, enabling the dissipation of detectable technosignatures. We emphasize that this transition is a qualitative implication requiring dedicated modeling of population and technosignature dynamics for quantitative validation.

The model's limitations are substantial and explicitly acknowledged. Calibration on twelve data points provides an existence proof, not a definitive parameter set. The binomial test for US turn accuracy does not support statistically significant predictive power. The compensation hypothesis remains a falsifiable conjecture. These limitations do not diminish the theoretical value of the framework; they define the agenda for its development.

Civilizations may not destroy themselves in nuclear fire. They may not exhaust their resources. They may simply lose the capacity for collective action at scale, their elites optimizing extraction over exploration, their technology turned inward for control rather than outward for expansion, until---at the final threshold---they fade from the observable universe entirely. The EDAP model provides a formal language for investigating this hypothesis and a computational platform for testing it against historical and future data.
